\documentclass[12pt, a4paper]{article}

\usepackage{ctex}
\usepackage{amsmath, amssymb, amsthm}
\usepackage{geometry}
\usepackage{fancyhdr}
\usepackage{enumitem}
\usepackage{xcolor}
\usepackage{booktabs}
\usepackage{hyperref}

\geometry{margin=2.4cm}
\pagestyle{fancy}
\fancyhf{}
\rhead{线性代数}
\lhead{学生姓名}
\rfoot{第 \thepage 页}
\renewcommand{\headrulewidth}{0.4pt}

\definecolor{okgreen}{RGB}{0,110,60}
\definecolor{warnred}{RGB}{170,30,30}
\definecolor{llmgray}{RGB}{90,90,90}

\newcommand{\problem}[1]{\subsection*{第 #1 题}}
\newcommand{\solution}{\textbf{解答。}}
\allowdisplaybreaks[3]

\begin{document}

\begin{center}
    {\LARGE\bfseries 作业解答} \\[0.5em]
    {\large 线性代数} \\[0.3em]
    学生姓名 \quad | \quad 2026-10-08
\end{center}

\hrule
\vspace{1em}
\setlength{\parskip}{0.35em}
\problem{Q1}

Can the graph of a function have more than one tangent at a given point? If so, graph your answer. If not, explain why.

\solution

In general, a smooth function has exactly one tangent line at each point.

However, a curve (not necessarily the graph of a function) can have multiple tangents at a self-intersection point.

For a function graph, the tangent is unique where the derivative exists.

\textbf{答案：}
\[\boxed{\text{Unique tangent where } f'(a) \text{ exists.}}\]

\vspace{0.8em}
\hrule
\vspace{0.8em}

\problem{Q2}

Is there a function whose graph doesn't have a tangent at some point? If so, graph your answer. If not, explain why.

\solution

Yes. A function may fail to have a tangent at a point where it is not differentiable.

Examples: $f(x) = |x|$ at $x = 0$ (corner), $f(x) = x^{1/3}$ at $x = 0$ (vertical tangent/cusp).

\textbf{答案：}
\[\boxed{\text{Yes: } f(x) = |x| \text{ at } x = 0 \text{ (corner)}}\]

\vspace{0.8em}
\hrule
\vspace{0.8em}

\problem{P1}

Suppose $\delta$ is a positive real number. How do you describe all real numbers $x$ that are within $\delta$ of 0?

\solution

x within δ of 0: $|x| < \delta$, i.e. $-\delta < x < \delta$

\textbf{答案：}
\[\boxed{$|x| < \delta$, i.e. $-\delta < x < \delta$}\]

\textbf{(a)}
The set of all $x$ within $\delta$ of 0: $-\delta < x < \delta$
$\boxed{-\delta < x < \delta}$

\textbf{(b)}
$|x| < \delta$
$\boxed{|x| < \delta}$

\textbf{(c)}
$(-\delta, \delta)$
$\boxed{(-\delta, \delta)}$

\vspace{0.8em}
\hrule
\vspace{0.8em}

\problem{P2}

Suppose $a$ is a real number and $\delta$ is a positive real number. How do you describe all real numbers $x$ that are within $\delta$ of $a$?

\solution

x within δ of a: $|x - a| < \delta$, i.e. $a - \delta < x < a + \delta$

\textbf{答案：}
\[\boxed{$|x - a| < \delta$, i.e. $a - \delta < x < a + \delta$}\]

\textbf{(a)}
[Number line with open circles at endpoints, center marked]
$\boxed{\text{See graph}}$

\textbf{(b)}
$a - \delta < x < a + \delta$
Equivalently: $-\delta < x - a < \delta$
$\boxed{a - \delta < x < a + \delta}$

\textbf{(c)}
$|x - a| < \delta$
$\boxed{|x - a| < \delta}$

\textbf{(d)}
$(a - \delta, a + \delta)$
$\boxed{(a - \delta, a + \delta)}$

\vspace{0.8em}
\hrule
\vspace{0.8em}

\problem{P3}

Suppose $a$ is a real number and $\delta$ is a positive real number. How do you describe all real numbers $x$ that are within $\delta$ of $a$, but not equal to $a$:

\solution

x within δ of a, not equal to a: $0 < |x - a| < \delta$

\textbf{答案：}
\[\boxed{$0 < |x - a| < \delta$}\]

\textbf{(a)}
[Number line with open circles at endpoints, center marked]
$\boxed{\text{See graph}}$

\textbf{(b)}
$a - \delta < x < a + \delta, \quad x \neq a$
Or equivalently: $a - \delta < x < a$ or $a < x < a + \delta$
$\boxed{a - \delta < x < a + \delta, \quad x \neq a}$

\textbf{(c)}
$0 < |x - a| < \delta$
$\boxed{0 < |x - a| < \delta}$

\textbf{(d)}
$(a - \delta, a) \cup (a, a + \delta)$
$\boxed{(a - \delta, a) \cup (a, a + \delta)}$

\vspace{0.8em}
\hrule
\vspace{0.8em}

\problem{P4}

Let $f(x) = x^2$.

\solution

Given: $f(x) = x^{2}$

\textbf{答案：}
\[\boxed{f(x) = x^{2}}\]

\textbf{(a)}
Need: $9.0 - 1 < x^{2} < 9.0 + 1$
i.e. $8.0 < x^{2} < 10.0$
Solving: $x \in (2.82842712474619, 3.16227766016838)$ (positive $x$ only)
$\boxed{(2.82842712474619, 3.16227766016838)}$

\textbf{(b)}
Need: $|x - 3| < \delta \Rightarrow |x^{2} - 9.0| < 1$
Boundaries: $x^{2} = 10.0 \Rightarrow x = \left[ -3.16227766016838, \  3.16227766016838\right]$
           $x^{2} = 8.0 \Rightarrow x = \left[ -2.82842712474619, \  2.82842712474619\right]$
Minimum distance from $x = 3$: $\delta = 0.16227766016838$
$\boxed{\delta = 0.16227766016838}$

\textbf{(c)}
\textit{\color{warnred} 未求解: 无法解析 ε-δ 计算子题}

\textbf{(d)}
Yes, this is true.
For any $\varepsilon > 0$, we can find $\delta > 0$ such that
$|x - a| < \delta \Rightarrow |x^{2} - L| < \varepsilon$.
This follows because $f(x)$ is continuous at $x = a$.
$\boxed{\text{Yes (continuity)}}$

\vspace{0.8em}
\hrule
\vspace{0.8em}

\problem{AP1}

Graph $y = 2x - x^2$. For which points $a$ is the tangent line to the curve at the point $(a, 2a - a^2)$ a horizontal line? How is the value of $2x - x^2$ at these points related to the values of $2x - x^2$ at other points?

\solution

Given: $y = - x^{2} + 2 x$

Derivative: $y' = 2 - 2 x$

Horizontal tangent where $y' = 0$: $2 - 2 x = 0$

Solutions: $x = \left[ 1\right]$

Points: $(1, 1)$

$f''(1) = -2 < 0$: this is a maximum.

The value $1$ is the maximum of $- x^{2} + 2 x$.

\textbf{答案：}
\[\boxed{(1, 1)}\]

\vspace{0.8em}
\hrule
\vspace{0.8em}

\problem{AP2}

The point $(2, 0)$ is on the curve $y = 2x - x^2$.

\solution

Given: $y = - x^{2} + 2 x$

Derivative: $y' = 2 - 2 x$

At point $(2, 0)$:

Slope: $f'(2) = -2$

Tangent line: $y = 4 - 2 x$

\textbf{答案：}
\[\boxed{4 - 2 x}\]

\textbf{(a)}
Draw line through $(2, 0)$ with slope $-2$
$\boxed{\text{Line with slope } -2}$

\textbf{(b)}
Slope at $x=2$: $-2$
Tangent: $y - 0 = -2(x - 2)$
$y = 4 - 2 x$
$\boxed{y = 4 - 2 x}$

\vspace{0.8em}
\hrule
\vspace{0.8em}

\begin{center}
\begin{tabular}{@{}lccc@{}}
\toprule
题号 & 引擎 & 已求解 & 回验通过项 \\
\midrule
    Q1 & - & 是 & 0/0 \\
    Q2 & - & 是 & 0/0 \\
    P1 & - & 是 & 0/0 \\
    P2 & - & 是 & 0/0 \\
    P3 & - & 是 & 0/0 \\
    P4 & - & 是 & 0/0 \\
    AP1 & - & 是 & 0/0 \\
    AP2 & - & 是 & 0/0 \\
\bottomrule
\end{tabular}
\end{center}

\noindent 共 8 题：已求解 8 题，其中通过符号回验 0 题。

\end{document}
